\documentclass[10pt,a4paper]{amsart}
\usepackage[margin=19mm]{geometry}
\usepackage{amsmath,amssymb,mathtools}
\usepackage[scr=rsfs,cal=euler]{mathalfa}
\usepackage{xfrac}
\usepackage[unicode,hidelinks,bookmarks=false]{hyperref}
\newcommand{\ie}{i.e.\ }
\newcommand{\cf}{cf.\ }
\newcommand{\resp}{respectively\ }
\newcommand{\Subsection}{\subsection}
\providecommand{\nobreakdash}{\nobreak\hbox{-}}
\setlength{\emergencystretch}{2em}
\multlinegap0pt
\usepackage{bm}
\usepackage{bbm}
\usepackage{dsfont}
\usepackage{xspace}
\usepackage{cancel}
\usepackage{mathtools}
\usepackage{amsthm}
\makeatletter
\@ifundefined{defi}{\newtheorem{defi}{Defi}}{}
\@ifundefined{Lemma}{\newtheorem{Lemma}[defi]{Lemma}}{}
\@ifundefined{prop}{\newtheorem{prop}[defi]{ Proposition}}{}
\makeatother
\title[Category O for p-adic rational Cherednik algebras]{Category O for p-adic rational Cherednik algebras}
\author{Fernando Pe\~na V\'azquez}
\date{Source: arXiv:2504.16699v2 (CC BY 4.0)}
\begin{document}
\maketitle

\noindent\emph{Source and licence.} Category O for p-adic rational Cherednik algebras. Version arXiv:2504.16699v2 (CC BY 4.0), distributed under the Creative Commons Attribution 4.0 International licence, as recorded in the arXiv licence metadata for this exact version and shown on the abstract page https://arxiv.org/abs/2504.16699v2. Full rights notice: licenses/arxiv-2504.16699-CC-BY-4.0.txt.\par\smallskip

\noindent\textit{Editorial scope.} Counted unit: the Prop whose source label is \texttt{prop ws decomposition of C} in the article (source0.tex lines 568--577, source label \texttt{prop ws decomposition of C}), delivered together with the article's own complete original proof of it (source0.tex lines 578--608, one \texttt{proof} environment of 31 lines). No mathematics is written, repaired or re-derived: statement and proof are the article's own text line for line. Required context retained uncounted: Lemma R admits a ws decomposition (lines 558--564). Every definition, display and label of those lines is retained unchanged. Omitted from this package and not redistributed: 1--557, 565--567, 609--1890. Nothing inside a retained excerpt is omitted or altered: every retained body is delivered exactly as the article's lines are written, so no omission receipt of any kind accompanies this package. Citations: the retained excerpts carry no citation command at all, so no bibliography entry of the article is transcribed and no reference is redistributed. Reference adaptation: none was needed and none was made; every reference inside the delivered unit resolves inside it, so no unresolved reference or citation is produced. Printed slips kept literal: no printed word, symbol, hypothesis or number of the counted statement or of its proof was altered. Layout and class: the article's own class is \texttt{amsart} with packages including babel, fontenc, tikz-cd, inputenc, fancyhdr, yhmath, mathrsfs, graphicx, wrapfig, caption, dsfont, ulem, enumitem, cutwin; the wrapper is the lane's pinned \texttt{amsart} editorial wrapper at 10pt on A4, which restates the theorem-like environments the retained text needs and adds the article's own macro definitions that the retained text needs (none), with extra wrapper packages (bm, bbm, dsfont, xspace, cancel, mathtools). The delivered numbering is the unit's own. Assets: no retained span contains a figure environment, a graphics-inclusion command, a PDF-page inclusion, a table or a TikZ picture; no image file is copied into this package, the package declares no asset dependency and no image file is redistributed with it. Licence: version arXiv:2504.16699v2 of this article is distributed under the Creative Commons Attribution 4.0 International licence, as recorded in the arXiv licence metadata for this exact version and shown on the abstract page https://arxiv.org/abs/2504.16699v2. A full rights notice is delivered at licenses/arxiv-2504.16699-CC-BY-4.0.txt. Characters: every retained excerpt is pure ASCII, so the delivered engine, XeLaTeX with its default Latin Modern text font, reports no missing character.
\begin{Lemma}\label{Lemma R admits a ws decomposition}
Choose any $a\otimes b \in \mathscr{C}$. Then we have:
\begin{equation*}
    \operatorname{ad}(\partial)(a\otimes b)=\operatorname{ad}(\partial)a\otimes b + a \otimes \operatorname{ad}(\partial)b.
\end{equation*}
In particular, the subspace  $\mathscr{C}\subset \mathscr{R}$ is closed under the action of $\operatorname{ad}(\partial)$.
\end{Lemma}

\begin{prop}\label{prop ws decomposition of C}
Assume that $\mathscr{A}$ and $\mathscr{B}$ satisfy the following condition:
\begin{itemize}
    \item Let $v=\sum_{\lambda\in I} v_{\lambda}$ with $\vert v \vert \leq \epsilon$. Then $\vert v_{\lambda}\vert \leq \epsilon$ for all $\lambda \in I$.
\end{itemize}
Then $\mathscr{C}$ admits a semi-simple weight space decomposition for $\operatorname{ad}(\partial)$ with $\Lambda(\mathscr{C})\subset\mathbb{Z}$. Furthermore, for any $n\in \mathbb{Z}$ we have:
\begin{equation*}
    \mathscr{C}^n=\widehat{\bigoplus}_{m\geq n}A_m\otimes B^{m-n}.
\end{equation*}
\end{prop}

\begin{proof}
First, by Lemma \ref{Lemma R admits a ws decomposition}, given $a\in A_n$, and $b\in B^m$ we have $a\otimes b \in \mathscr{C}^{n-m}$ for all $n,m\geq 0$. Let $v\in \mathscr{C}$. By construction of the completed tensor product of Banach spaces, there is a family $(a_i\otimes b_i)_{i\geq 0}\in \mathscr{A}\otimes_K\mathscr{B}$ satisfying the identities:
\begin{equation}\label{equation C admits a decomposition}
    v=\sum_{i\geq 0}a_i\otimes_Kb_i, \textnormal{  } \varinjlim_{i\geq 0}a_i\otimes_Kb_i=0.
\end{equation}
By definition of triangular decomposition, for each $i\geq 0$ there are unique identities:
\begin{equation*}
    a_i=\sum_{n\geq 0} a_i^n, \quad b_i=\sum_{m\geq 0} b_i^m,
\end{equation*}
such that $a_i^n\in A_n$, and $b_i^m\in B^m$ for all $n,m\geq 0$. As these are infinite sums:
\begin{equation}\label{equation 3 decomposition of C into ws}
    \varinjlim_n a_i^n=\varinjlim_n b_i^m=0.
\end{equation}
Hence, we may rearrange the terms in equation $(\ref{equation C admits a decomposition})$ to obtain:
\begin{equation*}
    v=\sum_{r\in \mathbb{Z}}\left(\sum_{i\geq 0}\sum_{r=n-m}a_i^n\otimes b_i^m  \right).
\end{equation*}
We are done if we are able to show that for each $r\in \mathbb{Z}$ the sum:
\begin{equation}\label{equation 2 C admits a ws decomposition}
    \sum_{i\geq 0}\sum_{r=n-m}a_i^n\otimes b_i^m,
\end{equation}
converges. It suffices to show that the principal term of $(\ref{equation 2 C admits a ws decomposition})$ converges to zero. As $\varinjlim_{i\geq 0}a_i\otimes_Kb_i=0$, for each $s\geq 0$ there is some $j\geq 0$ such that for all $i\geq j$ we have $\vert a_i\otimes_Kb_i \vert \leq \vert \pi \vert^{2s}$. In particular, we may assume that $\operatorname{max}(\vert a_i\vert, \vert b_i\vert)\leq \vert \pi \vert^s$ for all $i\geq j$. Thus, our assumptions on $\mathscr{A}$ and $\mathscr{B}$ imply that:
\begin{equation*}
    \operatorname{max}\{\vert a^n_i\vert, \vert b^m_i\vert\}_{i\geq j}^{n,m\geq 0}\leq \vert \pi \vert^s.
\end{equation*}
As there are finitely many $0\leq i \leq j$, it follows by $(\ref{equation 3 decomposition of C into ws})$ that there is $\overline{n}\geq 0$ such that:
\begin{equation*}
    \operatorname{max}\{\vert a^n_i\vert, \vert b^m_i\vert\}_{0\leq i\leq j}^{n,m\geq \overline{n}}\leq \vert \pi \vert^s.
\end{equation*}
Hence, the principal term of $(\ref{equation 2 C admits a ws decomposition})$ converges to zero, as we wanted to show.
\end{proof}

\end{document}
